%=====================================================================
\section{Lecture 6: Time-convolutionless equations and the NZ--TCL map}
\label{sec:lec6}

A time-nonlocal equation is exact but awkward: it needs the whole history of $\rhoS$, and
its truncations behave badly (\cref{fig:dynamics}). The time-convolutionless (TCL)
technique~\cite{Shibata1977,Chaturvedi1979,HanggiTalkner1977} trades the convolution for
a time-dependent generator, $\dot\rhoS=\Gc(t)\rhoS(t)$. This lecture derives it, computes it for
the qubit, shows how it is related to the NZ kernel, and discusses when it fails.

\subsection{Derivation}

Work in the interaction picture, $\dot{\tilde\rho}=\tilde\Lc_I(t)\tilde\rho$, and let
$\tilde G(t,s)$ be the full propagator from $s$ to $t$ (so $\tilde\rho(t)=\tilde G(t,s)\tilde\rho(s)$)
and $\tilde G(t,s)^{-1}=\tilde G(s,t)$ the backward propagator. The trick is to express the
\emph{past} in terms of the \emph{present}: $\tilde\rho(s)=\tilde G(s,t)\tilde\rho(t)$. Start from
the $\Qc$-equation solved as in \cref{sec:nz-derivation},
\begin{equation}
  \Qc\tilde\rho(t)=\hat G(t,0)\Qc\tilde\rho(0)+\int_0^t\!\dd s\,\hat G(t,s)\,\Qc\tilde\Lc_I(s)\,\Pc\,\tilde\rho(s),
  \qquad \hat G(t,s)=T\exp\Big[\int_s^t\!\Qc\tilde\Lc_I\Big],
\end{equation}
and substitute $\tilde\rho(s)=\tilde G(s,t)[\Pc+\Qc]\tilde\rho(t)$. Defining
\begin{equation}
  \Sigma(t)\equiv\int_0^t\!\dd s\,\hat G(t,s)\,\Qc\tilde\Lc_I(s)\,\Pc\,\tilde G(s,t),
\end{equation}
which is of order $\Lc_I$ and vanishes at $t=0$, we get $[1-\Sigma(t)]\Qc\tilde\rho(t)=\hat G(t,0)\Qc\tilde\rho(0)+\Sigma(t)\Pc\tilde\rho(t)$.
Wherever $1-\Sigma(t)$ is invertible (certainly for short times or weak coupling),
\begin{equation}
  \boxed{\;\partial_t\Pc\tilde\rho(t)=\tilde\Gc(t)\,\Pc\tilde\rho(t)+\tilde I(t),\qquad
  \tilde\Gc(t)=\Pc\tilde\Lc_I(t)\big[1-\Sigma(t)\big]^{-1}\Pc,\;}
  \label{eq:TCL}
\end{equation}
with $\tilde I(t)=\Pc\tilde\Lc_I(t)[1-\Sigma(t)]^{-1}\hat G(t,0)\Qc\tilde\rho(0)$, zero for factorized initial
states~\cite[Sec.~9.2]{BreuerPetruccione}. The equation is exact, time-local, and its
generator has a perturbative expansion obtained by expanding $[1-\Sigma]^{-1}=\sum_n\Sigma^n$:
\begin{align}
  \tilde\Gc^{(2)}(t)&=\int_0^t\!\dd t_1\,\Pc\tilde\Lc_I(t)\tilde\Lc_I(t_1)\Pc ,
  \label{eq:TCL2}\\
  \tilde\Gc^{(4)}(t)&=\int_0^t\!\dd t_1\!\int_0^{t_1}\!\dd t_2\!\int_0^{t_2}\!\dd t_3\,
  \Big[\avg{\tilde\Lc_I(t)\tilde\Lc_I(t_1)\tilde\Lc_I(t_2)\tilde\Lc_I(t_3)}
  -\avg{\tilde\Lc_I(t)\tilde\Lc_I(t_1)}\avg{\tilde\Lc_I(t_2)\tilde\Lc_I(t_3)}\notag\\
  &\hspace{3.2cm}-\avg{\tilde\Lc_I(t)\tilde\Lc_I(t_2)}\avg{\tilde\Lc_I(t_1)\tilde\Lc_I(t_3)}
  -\avg{\tilde\Lc_I(t)\tilde\Lc_I(t_3)}\avg{\tilde\Lc_I(t_1)\tilde\Lc_I(t_2)}\Big],
  \label{eq:TCL4}
\end{align}
where $\avg{\cdots}\equiv\Pc\cdots\Pc$ restricted to the system (the ``ordered cumulants'' of
van Kampen; see \cite[Sec.~9.2.4]{BreuerPetruccione} and
\cite{BreuerKapplerPetruccione2001}). Compare with the fourth-order NZ kernel: there, only
the \emph{reducible} pairing $(t\,t_1)(t_2t_3)$ is subtracted (by the $\Qc$'s), while here all
three products of pair contractions are subtracted, including the crossed and nested
ones. TCL4 is \emph{not} NZ4 with the convolution removed; it is a different
reorganization, which happens to contain NZ4 plus terms generated by the time-locality.

\subsection{TCL2 for the qubit: time-dependent rates}

In the Schr\"odinger picture, \cref{eq:TCL2} reads
\begin{equation}
  \dot\rhoS=\big[\Lc_S+\Gc^{(2)}(t)\big]\rhoS(t),\qquad
  \Gc^{(2)}(t)=\int_0^t\!\dd\tau\,K^{(2)}(\tau)\,\ee^{-\Lc_S\tau},
  \label{eq:TCL2-S}
\end{equation}
i.e.\ Redfield's generator \eqref{eq:redfield} with the upper limit at $t$ instead of
$\infty$. All the Markovian rates of Example~5.3 become running integrals:

\begin{exbox}[Example 6.1: the running rates]
Repeating the secular reduction of Example~5.3 with $\int_0^t$ instead of $\int_0^\infty$,
\begin{align}
  \Gamma_1(t)&=4\sin^2\theta\int_0^t\!\dd\tau\,S(\tau)\cos\omega_0\tau ,&
  \Gamma_1(t)\,r_z^{\rm eq}(t)&=-2\sin^2\theta\int_0^t\!\dd\tau\,\chi(\tau)\sin\omega_0\tau,\notag\\
  \Gamma_\varphi(t)&=4\cos^2\theta\int_0^t\!\dd\tau\,S(\tau),&
  \delta\omega(t)&=2\sin^2\theta\int_0^t\!\dd\tau\,S(\tau)\sin\omega_0\tau .
  \label{eq:running-rates}
\end{align}
Three features are general. (i) All rates vanish at $t=0$ as $t^1$ (since $K(0)\ne0$):
decay is initially quadratic in time, the Zeno regime. (ii) They approach their Markov
values on the bath correlation time $\tau_B$ ($=2/\gamma$ for the Lorentzian, $1/\omega_c$ or
$1/\pi T$ for the Ohmic bath). (iii) In between they can \emph{overshoot} and even become
negative: for the Lorentzian bath, $\Gamma_1(t)\propto\int_0^tS(\tau)\cos\omega_0\tau\,\dd\tau$ with
$S(\tau)\cos\omega_0\tau=\tfrac{g^2}2\ee^{-\gamma\tau/2}[1+\cos2\omega_0\tau]$ at resonance is monotonic, but
$\Gamma_\varphi(t)\propto\int_0^tS(\tau)\dd\tau=g^2\int_0^t\ee^{-\gamma\tau/2}\cos\Omega\tau\,\dd\tau$ oscillates and is
negative in windows when $\Omega>\gamma/2$. A negative rate in a time-local generator is the
standard signature of information flowing back from the bath~\cite{deVega2017}.
\end{exbox}

For pure dephasing ($\theta=0$) TCL2 is \emph{exact}: the generator on the coherence is
$-i\omega_0-\dot\Gamma(t)$ with $\Gamma$ of \cref{eq:dephasing-exact}, and $\dot\Gamma(t)=4\int_0^tS(\tau)\dd\tau=\Gamma_\varphi(t)$.
All higher ordered cumulants \eqref{eq:TCL4} vanish because the vertices commute and the
full fourth-order moment equals the sum of the three pair products. This is the cleanest
demonstration that ``order $n$'' means different things in different schemes: the same
model is second order in TCL and infinite order in NZ (Exercise~\ref{ex:5-4}).

\subsection{The NZ--TCL map}

Both equations describe the same map, so the kernel and the generator determine each
other. Writing the NZ equation for $\Pi(t)$ and multiplying by $\Pi(t)^{-1}$,
\begin{equation}
  \boxed{\;\Gc(t)=\Lc_S+\int_0^t\!\dd s\,K(t-s)\,\Pi(s)\,\Pi(t)^{-1},\qquad
  \Gc(t)=\dot\Pi(t)\Pi(t)^{-1},\;}
  \label{eq:TCLNZ}
\end{equation}
an exact relation wherever $\Pi(t)$ is invertible~\cite{ChruscinskiKossakowski2010}.
Conversely, given $\Gc$ one reconstructs $\Pi=T\exp\int\Gc$ and then $K$ from the Laplace
relation of Exercise~\ref{ex:5-4}. Nestmann, Bruch and Wegewijs~\cite{Nestmann2021} turned
\cref{eq:TCLNZ} into a fixed-point iteration, $\Gc\mapsto\Lc_S+\int K(t-s)\,T\ee^{\int_t^s\Gc}\dd s$, that
converges to the exact generator and makes explicit how a time-nonlocal kernel
``generates'' a time-local evolution: the generator at time $t$ is the kernel averaged over
the past with the propagator as weight. Inserting $K\approx K^{(2)}$ and $\Pi\approx\Pi_0$ gives
back $\Gc^{(2)}$; inserting $K^{(2)}+K^{(4)}$ with $\Pi$ to the required order reproduces
\cref{eq:TCL4} and exhibits the extra subtractions as the $\Pi^{(2)}\Pi_0^{-1}$ corrections
(Exercise~\ref{ex:6-3}). Reimer \emph{et al.}~\cite{Reimer2019} compare these and three further exact
formulations with attention to complete positivity and divisibility.

\subsection{When the generator does not exist}
\label{sec:tcl-singular}

$\Gc(t)$ is defined only where $\Pi(t)$ is invertible. For the transverse model
($\theta=\pi/2$) in the rotating-wave approximation and a Lorentzian bath at $T=0$, the
single-excitation sector is exactly solvable: the excited amplitude obeys
$\ddot c_e+\tfrac\gamma2\dot c_e+g^2c_e=0$ at resonance, so
\begin{equation}
  c_e(t)=\ee^{-\gamma t/4}\Big[\cosh\tfrac{dt}{4}+\tfrac{\gamma}{d}\sinh\tfrac{dt}{4}\Big],\qquad d=\sqrt{\gamma^2-16g^2},
\end{equation}
and $\rho_{ee}(t)=|c_e(t)|^2\rho_{ee}(0)$, $\rho_{eg}(t)=c_e(t)\rho_{eg}(0)$. The exact generator is of
GKSL form with the single time-dependent rate $\gamma(t)=-2\Re[\dot c_e/c_e]$ and shift
$-2\Im[\dot c_e/c_e]$~\cite[Sec.~10.1]{BreuerPetruccione,Vacchini2010}. For $16g^2>\gamma^2$ the
amplitude oscillates and crosses zero at finite times ($t\approx7.3/\omega_0$ for our $g=0.3$,
$\gamma=0.5$), where $\gamma(t)\to\infty$: the map is still perfectly well defined and
completely positive, but it is not invertible, and no time-local generator exists there.
The NZ kernel has no such problem. ``Local versus nonlocal'' is therefore a statement
about the representation, not about the physics~\cite{ChruscinskiKossakowski2010}.
Without the rotating-wave approximation the zero of $c_e$ is lifted (the counter-rotating
and $\sz$ terms feed the coherence from other components), but the near-singularity
survives as the large excursion of the exact generator in \cref{fig:kernels} (right) near
$t\omega_0\approx12$, where the smallest singular value of $\Pi(t)$ drops to $10^{-2}$.

\subsection{What to take from the numerical figures}

\Cref{fig:dynamics,fig:kernels} illustrate a hierarchy that holds quite generally for
Gaussian baths and weak-to-moderate coupling: Markov $<$ NZ2 $\lesssim$ TCL2 in
accuracy at short times, and Markov $\approx$ TCL2 $>$ NZ2 at long times, with NZ2 prone to
unphysical overshoots. The reason is structural: TCL2 is exact to second order in the
coupling \emph{for the generator}, and errors in the generator stay bounded; NZ2 is exact
to second order for the kernel but the convolution with the exact history amplifies the
missing fourth-order terms. Smirne and Vacchini~\cite{SmirneVacchini2010} give a careful
side-by-side for the two-level system; the lesson has been known since the early
comparisons of the two schemes in chemical physics (see the references in
\cite{BreuerKapplerPetruccione2001}).

\subsection{Exercises}

\begin{exercise}\label{ex:6-1}
Show that $\Sigma(0)=0$, that $\Sigma$ is of first order in $\Lc_I$, and derive
\cref{eq:TCL2,eq:TCL4} from \cref{eq:TCL}. Explain why odd orders vanish.
\end{exercise}

\begin{exercise}\label{ex:6-2}
Evaluate the running rates \eqref{eq:running-rates} in closed form for the Lorentzian bath
at $T=0$ and find the time windows in which $\Gamma_\varphi(t)<0$. Plot $\Gamma_\varphi(t)$, $\Gamma_1(t)$ and
$r_z^{\rm eq}(t)$ for the parameters of \cref{fig:dynamics} and compare with the script's output.
\end{exercise}

\begin{exercise}\label{ex:6-3}
Expand \cref{eq:TCLNZ} to fourth order using $\Pi=\Pi_0+\Pi^{(2)}+\dots$ and $K=K^{(2)}+K^{(4)}+\dots$,
and show that the extra terms relative to $\int K^{(4)}\Pi_0\Pi_0^{-1}$ are exactly the
crossed and nested subtractions in \cref{eq:TCL4}.
\end{exercise}

\begin{exercise}\label{ex:6-4}
For the rotating-wave model of \cref{sec:tcl-singular}, compute the exact NZ kernel
$\hat K_{ee}(z)$ for the population and show that it is finite and analytic for
$\Re z>-\gamma/4$ even when $c_e$ has zeros. Where are the singularities of
$\hat\Gc$, if one tries to define it by Laplace transform?
\end{exercise}

\clearpage
